\origpage{348}{N. Robertson, P.D. Seymour / Journal of Combinatorial Theory, Series B 90 (2004) 325--385}
The three separations $(H, H')$, $(A \cap H', B \cup H)$ and $(A \cup H, B \cap H')$ all have $\lambda$-order at most that of $(A, B)$, because $H \cap H' \subseteq A \cap B$ and every edge of $\Gamma^*$ split by one of these separations is split by $(A, B)$.

If $H \subseteq A$ then $A \cup H' = \Gamma$, and so $(H', H) \notin \mathscr{T}$ since $(A, B) \in \mathscr{T}$; while if $H \subseteq B$ then $B \cup H = \Gamma$, and so $(H', H) \notin \mathscr{T}'$ since $(B, A) \in \mathscr{T}'$. In either case it follows that $(H', H)$ does not belong to both $\mathscr{T}, \mathscr{T}'$, and so $(H, H')$ belongs to at least one of them.

Suppose that $(H, H')$ belongs to both of $\mathscr{T}, \mathscr{T}'$. Now $(A \cap H', B \cup H) \notin \mathscr{T}'$, from the second tangle axiom, since $(H, H'), (B, A) \in \mathscr{T}'$ and $H \cup B \cup (A \cap H') = \Gamma$. Consequently $(B \cup H, A \cap H') \in \mathscr{T}'$. But $(A \cap H', B \cup H) \in \mathscr{T}$, since $(A, B) \in \mathscr{T}$; and since $(A \cap H', B \cup H)$ has $\lambda$-order at most that of $(A, B)$, and $(A, B)$ is the $(\mathscr{T}, \mathscr{T}')$-distinction, it follows that equality holds, and so $(A \cap H', B \cup H) = (A, B)$, that is, $A \subseteq H'$ and $H \subseteq B$. Similarly, $(B \cap H', A \cup H) \notin \mathscr{T}$, since $(H, H'), (A, B) \in \mathscr{T}$, and $H \cup A \cup (B \cap H') = \Gamma$. Consequently $(A \cup H, B \cap H') \in \mathscr{T}$. But $(B \cap H', A \cup H) \in \mathscr{T}'$, since $(B, A) \in \mathscr{T}'$, and so this separation has the same $\lambda$-order as $(A, B)$, and therefore $(A \cup H, B \cap H') = (A, B)$, that is, $H \subseteq A$ and $B \subseteq H'$. But we already showed that $A \subseteq H'$ and $H \subseteq B$, and so $E(H) = \emptyset$. Since $\Delta$ includes a region of $\Gamma^*$ and hence an edge of $\Gamma$, this is impossible.

It follows that $(H, H')$ belongs to exactly one of $\mathscr{T}, \mathscr{T}'$. Since its $\lambda$-order is at most that of $(A, B)$, and $(A, B)$ is the $(\mathscr{T}, \mathscr{T}')$-distinction, it follows that equality holds, and so $(H, H') = (A, B)$ or $(B, A)$. The first is the desired result. If the second holds, then since $(H, H')$ has order less than $rep(\mathscr{T})$ and $(H, H') \notin \mathscr{T}$, it follows that there is a disc bounded by $F$ different from $\Delta$; and so $\Delta'$ is a disc. Since in this case $\Delta' = ins(F)$, we deduce that again the desired result holds. $\square$

Given a tangle $\mathscr{T}$ in $\Gamma$, a circuit $F$ of $\Gamma^*$ is a $\mathscr{T}$-enclave if $\frac{1}{2}|E(F)| < rep(\mathscr{T})$, and there is a tangle $\mathscr{T}'$ in $\Gamma$ for which $(\Gamma \cap ins(F), \Gamma \cap out(F))$ is the $(\mathscr{T}, \mathscr{T}')$-distinction. We also call $F$ a $\mathscr{T}$-enclave around $\mathscr{T}'$. If $K$ is a subgraph of $\Gamma^*$, $\lambda(K)$ denotes $\sum_{e \in E(K)} \lambda(e)$.

\textbf{8.5.} \textit{Let $\mathscr{T}$ be a tangle in $\Gamma$, and let $F$ be a $\mathscr{T}$-enclave around $\mathscr{T}'$, and let $u, v \in V(F)$ be distinct. Let $F_1, F_2$ be the two paths of $F$ between $u$ and $v$, and let $P$ be a path of $\Gamma^*$ between $u$ and $v$ with no other vertex or edge in common with $F$. Suppose that $|E(P)| \le |E(F_1)|$, and if equality holds then $\lambda(P) \le \lambda(F_1)$. Then}

\hspace*{12pt}$\bullet$ \textit{$(\Gamma \cap ins(P \cup F_2), \Gamma \cap out(P \cup F_2)) \in \mathscr{T}'$,}

\hspace*{12pt}$\bullet$ \textit{if $U(P) \nsubseteq ins(F)$ then $ins(P \cup F_2) \cap ins(F) = U(F_2)$,}

\hspace*{12pt}$\bullet$ \textit{$|E(P)| \ge |E(F_2)|$, and if equality holds then $\lambda(P) > \lambda(F_2)$.}

\textbf{Proof.} Let $C_i$ be the circuit $P \cup F_i$ ($i = 1, 2$). Certainly $|E(C_2)| \le |E(F)|$, and if equality holds then $\lambda(C_2) < \lambda(F)$, because the $\lambda(e)$'s are rationally independent. Since $F$ is a $\mathscr{T}$-enclave around $\mathscr{T}'$, it follows that $(\Gamma \cap ins(C_2), \Gamma \cap out(C_2)) \notin \mathscr{T} \setminus \mathscr{T}'$. Hence the first assertion of the theorem holds. It follows that $ins(F) \nsubseteq ins(C_2)$, because $(\Gamma \cap out(F), \Gamma \cap ins(F)) \in \mathscr{T}'$, and so the second assertion holds. Suppose that the
